\section{Design}

\subsection{Overview}
\begin{figure*}[t]
	\centering
	\includegraphics[width=0.7\textwidth]{figure/overall_design.png}
	\caption{Overall Architecture of Our Approach}
	\label{overall_design}
\end{figure*}

We propose a hybrid framework that combines offline static analysis with runtime dynamic monitoring for rule conflict detection in the smart home. As shown in Figure~\ref{overall_design}, the framework consists of two phases. In the static phase, the system parses automation configurations together with the Physical Channel Configuration to build enhanced rule models, applies formal verification to construct the RIDG, and then uses the Entity Safety Configuration to prune the RIDG and extract the RCDG, which retains the interaction edges corresponding to candidate rule conflicts. In the runtime phase, the system monitors rule events, checks whether triggered rules correspond to nodes in the RCDG, and uses assertion verification with the current system state to determine whether a candidate conflict edge is activated. If so, the candidate rule conflict associated with that edge is reported as an impending rule conflict. In essence, the static phase first determines whether a rule interaction is physically relevant in the current deployment and then checks whether it violates the safety boundary of the target application scenario; the runtime phase confirms whether such a candidate conflict is actually activated.

\subsection{Static Analysis}

We begin by clarifying how automation rules interact. The action of one rule may affect the trigger, condition, or action of another rule either directly or through physical channels. We therefore consider six types of rule interaction: Trigger Interaction, Condition Interaction, Action Interaction, Indirect Trigger Interaction, Indirect Condition Interaction, and Indirect Action Interaction. These categories provide the basis for constructing the RIDG, after which the Entity Safety Configuration is applied to the identified interaction edges to retain candidate rule conflicts.

Automation rules are usually represented by Trigger (T), Condition (C), and Action (A). However, the TCA model alone cannot represent indirect interaction through physical channels. We therefore extend it to the TCAE model by introducing environmental attributes $E_T$, $E_C$, and $E_A$, where each element is written as $e=(zone, channel, trend)$. Here, $zone$ denotes the affected zone, $channel$ denotes the physical quantity, and $trend$ denotes whether the quantity increases or decreases. For example, if a rule turns on a bedroom heater, its action may include $\langle bedroom, temperature, increase \rangle$ in $E_A$. A rule is thus modeled as $R=\langle T, C, A, E_T, E_C, E_A \rangle$, which allows the system to represent indirect interaction and zone boundaries. The Physical Channel Configuration is used to determine these environmental attributes from the deployment setting.

After modeling automation rules with TCAE, we use formal verification to identify pairwise rule interaction and organize the results as a directed graph. We define the RIDG as $G_I=(V,E_I)$, where each vertex in $V$ denotes an automation rule and each directed edge in $E_I$ denotes that one rule may affect another through a direct or indirect interaction. For each pair of rules $R_i$ and $R_j$, the verifier checks whether they satisfy any formal condition in Table~\ref{Formal_Analysis_Expression}. If so, a directed edge $R_i \rightarrow R_j$ is added to the RIDG. The RIDG therefore captures the general rule interactions that are logically or physically possible under the current rule set and deployment setting.

Table~\ref{Formal_Analysis_Expression} summarizes the formal verification rules. Assume $R_1$ and $R_2$ are configured in the same system. The symbol $\rightarrow$ denotes that an action or environmental effect facilitates a trigger or condition, whereas $\nrightarrow$ denotes that it prevents a condition from being satisfied. For physical-channel attributes, $\bot$ denotes that two attributes have the same $zone$ and $channel$ but opposite $trend$, or that two actions conflict with each other. The symbol $=$ denotes attribute equivalence. Here, $E_{A_1}$ denotes the set of environmental effects induced by the actions in $A_1$, $E_a$ denotes the set of environmental effects induced by a specific action $a$, and $D_a$ denotes the target device of action $a$. When a relation is written over a set of actions or environmental effects, it is interpreted existentially. For example, $A_1 \rightarrow C_2$ means that there exists $a_1 \in A_1$ such that $a_1 \rightarrow C_2$, and $E_{A_1} \perp E_{C_2}$ means that there exist $e_1 \in E_{A_1}$ and $e_2 \in E_{C_2}$ such that $e_1 \perp e_2$.

\begin{table*}[t]
	\centering
	\small
	\caption{Formal Verification Rules for Rule Interaction}
	\label{Formal_Analysis_Expression}
	\begin{adjustbox}{width=0.95\textwidth}
		\begin{tabular}{c|c}
			\hline
			\textbf{Category} & \textbf{Expression}\\
			\hline
			\text{Trigger Interaction} & $(\exists a_{1}\in A_{1})\wedge(a_{1}\rightarrow T_{2})$ \\
			\hline
			\text{Condition Interaction} & $\left((\exists a_{1}\in A_{1}: a_{1}\nrightarrow C_{2})\wedge(R_{1}\neq R_{2})\right)\vee\left((\exists a_{1}\in A_{1}: a_{1}\rightarrow C_{2})\wedge(R_{1}\neq R_{2})\right)$ \\
			\hline
			\text{Action Interaction} & $(\exists a_{1}\in A_{1})\wedge(\exists a_{2}\in A_{2})\wedge(a_{1}\perp a_{2})$ \\
			\hline
			\text{Indirect Trigger Interaction} & $(\exists a_{1}\in A_{1})\wedge(\exists e_{1}\in E_{a_{1}})\wedge(\exists e_{2}\in E_{T_{2}})\wedge(e_{1}=e_{2})$ \\
			\hline
			\text{Indirect Condition Interaction} & $(\exists e_{1}\in E_{A_{1}})\wedge(\exists e_{2}\in E_{C_{2}})\wedge\left((e_{1}\perp e_{2})\vee(e_{1}=e_{2})\right)\wedge(R_{1}\neq R_{2})$ \\
			\hline
			\text{Indirect Action Interaction} & $(\exists a_{1}\in A_{1})\wedge(\exists a_{2}\in A_{2})\wedge(D_{a_{1}}\neq D_{a_{2}})\wedge(\exists e_{1}\in E_{a_{1}})\wedge(\exists e_{2}\in E_{a_{2}})\wedge(e_{1}\perp e_{2})$ \\
			\hline
		\end{tabular}
	\end{adjustbox}
\end{table*}

Not all rule interactions should be retained as candidate rule conflicts. We therefore use the Entity Safety Configuration to prune the RIDG and extract the RCDG. Formally, the RCDG is a subgraph of the RIDG, denoted by $G_C=(V_C,E_C)$, where $V_C \subseteq V$ and $E_C \subseteq E_I$. Each edge in $E_C$ represents a rule interaction that may move a safety-sensitive entity away from its preferred safe state and therefore requires further runtime checking. The Entity Safety Configuration defines preferred safe states for safety-sensitive entities, such as preferring a door lock to remain locked and a sprinkler to remain enabled. These settings can be initialized from common IoT safety knowledge and further refined by the user.

To construct the RCDG, we assign a safety score ($sf$) to each interaction edge in the RIDG. The score is increased when the resulting state is consistent with the preferred safe state and decreased otherwise; the default value is zero. An edge is retained in the RCDG if its interaction result moves a safety-sensitive entity away from its preferred safe state, for example when $sf<0$. Otherwise, the edge is pruned. The safety score serves as the core pruning signal, while edge retention is interpreted according to the semantics of the corresponding interaction category; for example, action-based conflicts additionally require contradictory effects on a safety-sensitive entity or physical channel. The static phase also records the conflict type and a translation template for each retained edge for later reporting.

\subsection{Runtime Monitoring}

Runtime dynamic monitoring focuses on rule events and their corresponding nodes in the RCDG. When an automation rule $R_i$ is triggered, the system first checks whether $R_i$ has a corresponding node in the RCDG. If not, the current event does not require further conflict checking. If such a node exists, the system further checks whether it participates in any outgoing candidate conflict edges. If no such edge exists, no additional runtime verification is required. Otherwise, the system selects the corresponding candidate conflict edges for subsequent assertion verification. The runtime context is updated when the affected device state is later changed by user intervention or by another automation rule, so that expired rule effects are removed from subsequent checking.

Even when an edge $R_i \rightarrow R_j$ exists in the RCDG, the corresponding candidate conflict edge is not necessarily activated at runtime. Assertion verification therefore checks whether the candidate conflict edge indicated by static analysis is supported by runtime evidence. We use $obs(X)$ to denote that event $X$ is observed, including trigger events $obs(T)$, successful condition checks $obs(C)$, failed condition checks $obs(\neg C)$, and action execution $obs(A)$. We use $intime(X,Y,\delta)$ to denote that the time interval between $X$ and $Y$ is smaller than $\delta$, where the default threshold is $\delta = 0.01\,\mathrm{s}$. Because direct and indirect candidate rule conflicts share the same runtime evidence pattern, we use four assertion templates, summarized in Table~\ref{Assertion_Verification}. For Action-Based Conflict, the framework reports an impending conflict before the second action is executed: observing $T_2$ and $C_2$ means that the second rule has entered the action-ready state, and no additional $intime(\cdot)$ constraint is required because the conflict may arise even when the two actions are separated in time, as long as the first rule's effect remains valid in the current runtime context.

\begin{table}[t]
	\centering
	\small
	\caption{Runtime Assertion Verification Rules}
	\label{Assertion_Verification}
	\renewcommand{\arraystretch}{1.1}
	\begin{adjustbox}{width=\textwidth}
		\begin{tabular}{>{\centering\arraybackslash}m{4.2cm}|>{\raggedright\arraybackslash}p{9.2cm}}
			\hline
			\textbf{Category} & \multicolumn{1}{c}{\textbf{Verification Expression and Explanation}}\\
			\hline
			\makecell[c]{Trigger-Based Conflict}
			&
			\makecell[l]{
				$obs(T_1), obs(C_1), obs(A_1)$ \rcmtinline{$R_1$ is executed} \\
				$obs(T_2), obs(C_2), intime(A_1, T_2, \delta)$ \\
				\rcmtbelow{within $\delta$, $R_2$ is triggered and its condition is satisfied}
			}\\
			\hline
			\makecell[c]{
				Condition-Based Conflict \\
				(Make Conditions Forbidden)
			}
			&
			\makecell[l]{
				$obs(C_2)$ \rcmtinline{$C_2$ is initially satisfied} \\
				$obs(T_1), obs(C_1), obs(A_1)$ \rcmtinline{$R_1$ is executed} \\
				$obs(\neg C_2), intime(A_1, \neg C_2, \delta)$ 
				\rcmtinline{within $\delta$, $C_2$ is unsatisfied}
			}\\
			\hline
			\makecell[c]{
				Condition-Based Conflict \\
				(Make Conditions Satisfied)
			}
			&
			\makecell[l]{
				$obs(\neg C_2)$ \rcmtinline{$C_2$ is initially unsatisfied} \\
				$obs(T_1), obs(C_1), obs(A_1)$ \rcmtinline{$R_1$ is executed} \\
				$obs(C_2), intime(A_1, C_2, \delta)$
				\rcmtinline{within $\delta$, $C_2$ is satisfied}
			}\\
			\hline
			\makecell[c]{Action-Based Conflict}
			&
			\makecell[l]{
				$obs(T_1), obs(C_1), obs(A_1)$ \rcmtinline{$R_1$ is executed} \\
				$obs(T_2), obs(C_2)$ 
				\rcmtinline{$R_2$ reaches the action-ready state}
			}\\
			\hline
		\end{tabular}
	\end{adjustbox}
\end{table}

When all predicates in the corresponding row are satisfied, the system treats the candidate conflict edge as activated and reports the associated candidate rule conflict as an impending rule conflict. The runtime controller collects the required runtime evidence and generates a detection report that combines the static conflict characterization, such as conflict type, physical channels, and safety score ($sf$), with the temporal evidence observed at runtime.